\chapter{微调方案怎样选：训练目标与更新参数是两回事}

\section{先分清两个维度}

“用什么误差教模型”和“允许哪些参数改变”是两个独立选择。SFT、偏好学习、强化学习描述学习目标和反馈形式；训练分类头、LoRA、QLoRA、局部解冻、全量更新描述参数与存储方式。不能把它们排成互相替代的一列。例如当前方案同时属于 SFT 和 QLoRA；DPO 也可以只更新 LoRA，而不必更新全部底座。

另一条独立维度是输出形式：可以生成 JSON 文本，可以直接输出技能分类概率，也可以预测连续动作。它们不由“使用 VLM”这三个字自动决定。小鸭子已有可靠的底层技能时，首先要学的是\textbf{根据观察选择技能、参数和结束时机}；是否训练关节控制器是另一项工作。

\begin{longtable}{@{}>{\raggedright\arraybackslash}p{2.0cm}>{\raggedright\arraybackslash}p{3.0cm}>{\raggedright\arraybackslash}p{5.0cm}>{\raggedright\arraybackslash}p{4.5cm}@{}}
\hline 方式 & 更新什么 & 输入、输出与数据 & 成本与能力边界 \\
\hline
\endfirsthead
\hline 方式 & 更新什么 & 输入、输出与数据 & 成本与能力边界 \\ \hline
\endhead
仅提示词 & 无参数更新 & 图像、指令、状态进入原模型；提示其生成技能 JSON；可放少量示例。 & 无训练成本但仍有推理成本；接口描述会进入上下文，不能保证遵从。\\
冻结骨干加头 & 新增 linear 或 MLP 头 & 先获得联合图文状态特征，再监督技能及参数；输入需有正确标签。 & 头很小，特征可缓存；适应能力受冻结特征限制，任意文本回复还需额外路径。\\
局部解冻 & 投影器或选定后层 & 可用 JSON 示范或分类标签，目标与输出由设计决定。 & 比全量少更新，但仍有前后向开销；要防止破坏原有对齐。\\
LoRA & 各层少量低秩矩阵 & 可用于文本 SFT，也可与动作头联合训练。 & 优化器存储较小，骨干通常仍占较多显存；能力需实测。\\
QLoRA & 量化冻结底座旁的 LoRA & 当前为双图、指令、状态到下一步 JSON。 & 已验证本机可运行；量化减少存储，并非免除骨干计算。\\
全量 SFT & 所选模型全部参数 & 通常需较多且多样的高质量示范；输出仍可选择。 & 梯度、优化器和激活需求大，6 GB 本机未验证；易过拟合或遗忘。\\
DPO 等偏好学习 & 可用 LoRA 或全量更新 & 同一输入配更好和更差的输出，学习相对偏好。 & 需要可信偏好对；不是已有单一正确标签直接换个开关。\\
环境强化学习 & 上层策略或底层控制器 & 状态、动作、轨迹、奖励和终止信号；目标是累计回报。 & 需要大量交互与奖励设计；能优化长期任务，但训练复杂度更高。\\
\hline
\end{longtable}

表中只有“提示词驱动现有技能”和“当前 JSON QLoRA 训练链路”已经在本项目中实现。分类头、投影器微调、全量 SFT、DPO 与上层强化学习均是可研究方案，不是已经配置好的启动命令。它们的准确率和训练时间不能从当前 QLoRA 日志推算。

\section{方案一：冻结 VLM，只训练技能预测头}

这通常是一个很值得做的对照实验。先让冻结的 VLM 将允许输入变成向量，再让一个小网络学习如何将这个向量映射到技能。冻结特征后训练新层、或预先保存特征再训练小模型，是标准迁移学习流程；官方 Keras 指南也明确说明了缓存能减少重复骨干计算，以及动态数据增强会使固定缓存不再适用。\cite{transfer-feature-cache}

\subsection{必须缓存联合输入特征，不能偷看目标}

对于第 $i$ 条样本，先构造：
\[
z_i=F_{\theta_0}(I_i^{\mathrm{front}},I_i^{\mathrm{near}},
u_i,s_i,h_i),\qquad \theta_0\ \text{保持冻结}.
\]
这里 $u_i$ 为用户指令，$s_i$ 为允许状态，$h_i$ 为历史。对于因果 VLM，一种候选读出位置是\textbf{助手开始标记之前或完整输入末端的隐藏状态}，它已能注意到前面的图像、指令和状态；也可以比较输入位置池化。究竟哪种读出更适合技能预测需要验证，不应把“最后一个向量”当成已证明的最佳表示。

缓存时绝不能将目标 JSON 接在后面，再从回答末端取向量。那样特征已经读到了正确技能，训练和离线测试可能看起来极好，上线时却拿不到这样的特征。教师动作、规则阶段、仿真世界坐标、任务编号、文件名中的标签也不能偷偷混进输入。相同画面配“踢红球”和“踢蓝球”时，联合特征应允许因指令不同而不同；只缓存图像向量后直接分类，会丢掉语言条件。

缓存还应固定底座版本、提示词、处理器、取特征位置和推理模式，并为缓存记录输入与配置哈希。图像或文字增强改变了输入，就应重新提取对应特征；一旦开始更新骨干或 LoRA，旧联合特征也会失效。数据的 train、val、test 划分必须在特征训练前固定，标准化统计只从训练部分估计。

若每条只保存一个 2048 维 float16 向量，9500 条训练样本的数值本体约为 $9500\times2048\times2$ 字节，即 37.1 MiB，另有标签和元数据。这是\textbf{所提缓存格式的体积估算}，不是已经产出的文件。离线提取仍须为每条运行一次 VLM；实时机器人每获得新画面，也仍要运行一次骨干。缓存主要加快反复训练小头，不能把未来从未见过的画面提前算好。

\subsection{四个头分别预测什么}

可以在共享特征上建立一个 linear 头，或先接小型 MLP，再输出四组 logits：
\[
(a^{\mathrm{type}},a^{\mathrm{skill}},a^{\mathrm{duration}},a^{\mathrm{color}})
=H_\phi(z).
\]
一个与当前接口对应的候选设计如下。
\begin{longtable}{@{}>{\raggedright\arraybackslash}p{3.0cm}>{\raggedright\arraybackslash}p{5.0cm}>{\raggedright\arraybackslash}p{6.8cm}@{}}
\hline 预测头 & 类别 & 使用方法 \\
\hline
type & 4 类 & skill、answer、clarify、finish，所有样本监督。\\
skill & 9 类 & stand、前后左右移动、左右转向、左右踢；仅动作样本计算该项损失。非动作输出 null。\\
duration & 51 个候选格点 & $0,0.02,\ldots,1.00$ 秒；根据 type 和 skill 屏蔽不合法值。\\
color & 4 类 & red、yellow、blue、null；所有样本可监督，学习无颜色目标的情况。\\
\hline
\end{longtable}

这个设计不能随意把任何头的最大概率组合起来。例如非动作必须 \texttt{skill=null, seconds=0}；移动时长只能为 0.02 至 0.60 秒，站立和转向为 0.02 至 1.00 秒；两种踢球固定 0.50 秒。可以先选择 type，再在对应合法技能与时长集合中解码。训练 duration 项也使用正确技能对应的合法集合，推理则必须使用预测技能，不能继续用真值来替模型纠错。

颜色不能简单规定“非动作一律 null”。本地示范中，完成踢红球任务的 finish 仍可保留 red；而回答视野问题通常为 null。哪些结构组合不合法可由接口确定，但“这次应该选蓝色还是红色”是需要学习的内容，不能用测试标签建立屏蔽表。

一种起点损失为：
\begin{align*}
\mathcal{L}_{\mathrm{head}}={}&
\lambda_t\operatorname{CE}(a^{\mathrm{type}},t)
+m_s\lambda_s\operatorname{CE}(a^{\mathrm{skill}},k)\\
&+m_d\lambda_d\operatorname{CE}(a^{\mathrm{duration}},d)
+\lambda_c\operatorname{CE}(a^{\mathrm{color}},c).
\end{align*}
其中 $m_s$ 只在 type 为 skill 时取 1；$m_d$ 只在需要学习可变时长的动作上取 1。固定踢球时长可直接由接口赋值，避免训练一个没有选择余地的分类项。各权重可先设为 1，再依据验证集调整，不能为了提高测试分数反复改测试配置。若少数技能严重不足，可以在训练集内采用类别权重或分层采样，同时保持评估分布清楚可追溯。

例如“向前一点”的正确标签可以为 skill、move\_forward、0.32、null。若模型把类型判成 finish，即使另一个头恰好选对 move\_forward，最后动作仍然错。评估必须查看联合结果，而不只是分别挑最好看的单头准确率。分类概率也不天然是可信置信度；若使用低置信度转为 clarify，需要在验证集上校准并检查是否滥用澄清来逃避任务。

\subsection{技能头怎样保留“你看到了什么”}

分类头只输出有限类别和参数，不能凭空生成自由中文 \texttt{reply}。一种完整部署方案是：动作、结束和固定澄清由头决定并填充模板；当 type 为 answer 时，再调用冻结 VLM 的文本生成能力回答当前画面。另一种方案是另训练回复解码路径，但那已超出“只训练小分类头”。

因此比较两种方案时，应分别测动作决策和视觉问答。不能把只会输出 skill ID 的小头称为已经拥有完整对话功能，也不能假定冻结底座的回答在新的仿真视角上一定正确。无论哪种输出方式，最终都要转换为同一个技能协议，再经过相同的时长约束、踢球条件和执行反馈闭环。

\section{其他参数更新方法何时有意义}

\subsection{投影器或局部解冻}

当图像特征对新相机外观、视角或任务细节不够适配时，可以研究更新视觉到语言的投影模块、语言末端若干层，或小部分视觉层。它比全冻结有更多调整空间，但“解冻一小块”并不自动意味着显存很小：只要梯度要穿过相关路径，就需要保存或重算激活。当前视觉编码器和投影器全部冻结，尚无这一方案的本机性能结论。

应先检查错误来自哪里。若颜色识别已对，错误集中在何时停止，先改善上层示范和时序监督通常比改视觉模块更直接。若目标在图中只有几个像素，换训练参数也未必能弥补输入分辨率缺失；需要一起评估图像设置、相机覆盖和数据质量。

\subsection{LoRA、QLoRA 与全量 SFT}

LoRA 用低秩矩阵表达权重变化，QLoRA 则把它与量化冻结底座结合。选择 $r$ 是容量与成本的权衡，不是准确率旋钮；更大 $r$ 需要更多可训练参数，也可能更容易记住有限示范。本次 $r=8$ 是经过显存验证的起点，不代表已经搜索过全部最优配置。\cite{lora-paper,qlora-paper}

全量 SFT 允许原模型大量参数直接更新，适应空间更大，但也需要存放更多梯度与优化器状态，且可能损伤原有视觉语言能力。“全量微调”仍然从预训练权重出发，和随机初始化后从零训练不同。现有 6 GB 环境没有验证全量方案；不能直接套用当前 QLoRA 的 15--17 小时估算。

如果冻结特征小头明显欠拟合、数据质量可靠，再考虑增大头、训练部分适配器或局部解冻，可以更清楚地识别新增复杂度是否有效。每改变一次可训练范围，都应单独记录数据、参数量、资源和评测结果，而不是同时改变所有设置后把改进归因于某一个名词。

\section{偏好学习和强化学习不是“更高级的微调按钮”}

DPO 使用同一输入下的偏好输出对，促使模型提高更好输出相对于更差输出的倾向。它需要明确的 chosen/rejected 数据，方法本身不等于让机器人在线探索。原论文将其作为不必显式训练奖励模型的偏好优化方式。\cite{dpo-paper}

对小鸭子，可以研究“目标偏左时，小幅左转”优于“直接踢球”的偏好，但不能仅因某个输出来自现有规则就自动认为它永远更好。有效偏好应结合执行结果、对齐程度和任务完成情况。当前数据主要是输入到单一教师动作，尚没有经过验证的偏好对或 DPO 训练代码。

强化学习则让策略通过环境交互，优化整个轨迹的累计奖励，例如正确球被踢动的奖励、错误球接触的惩罚、跌倒或超时惩罚。奖励和终止条件定义不好，策略可能找到“看起来得分但任务没完成”的行为。语言模型领域的 RLHF 展示了示范、偏好与后续强化优化可以组合，但这不直接证明同一流程能迁移到本机器人。\cite{instructgpt-paper}

还需区分层级：训练上层“选择左转还是踢球”的策略，可以继续复用底层控制器；训练底层步态或踢腿，是另一套状态、动作频率和奖励。本次上层 QLoRA 只模仿示范，没有在线奖励优化，也不会在每次日常执行动作后自动再训练一轮。

\section{怎样公平比较小头与当前 JSON 方案}

比较应使用相同的训练、验证、测试划分和独立场景种子，输入相同两幅图、同一指令及允许状态，采用同样的技能库与执行器。两种方法可以选不同输出表示，但不能给一方额外的目标坐标或未来反馈。任何特征缓存都只能由该条允许输入产生。

离线部分至少报告合法输出比例、动作样本上的九技能准确率和混淆矩阵、每种颜色准确率、时长误差与精确一致率、四字段联合准确率，以及视觉回答是否符合图像。分类头天然容易保证结构合法，JSON 模型需要生成正确结构；应分别报告原始输出质量和进入同一执行 guard 后的结果，避免比较口径不一致。

闭环部分使用相同新场景和任务说法，测正确颜色球的实际踢动成功率、错误球动作、搜索步数、耗时、失败与拒绝原因。还应测试目标最初在盲区、靠近时暂时出视野、执行被拒绝后的调整，以及“往前一点”完成后能否停止。离线标签存在多种合理解时，精确匹配低不一定说明任务差；反过来，离线匹配高也可能在连续执行中积累误差。

成本比较要包括特征提取时间、后续训练时间、峰值显存、实时每次决策延迟和总实验迭代成本。缓存方案不能只报小头训练几秒而省略首次 VLM 提取；当前 JSON 方案也不能只报权重小而忽略逐 token 生成时间。测试集用于最终报告，不用于反复选择头大小、阈值或提示词。

截至本报告快照，已经具备的是可恢复的 QLoRA 训练、独立数据划分和初始离线基线。冻结联合特征加多头分类器是一项清晰、可复用现有数据的后续对照方案；其收益必须通过上述共同标准验证，不能先把它写成更快且更准的既定结论。
